% Pursuing Stacks (cotes 134-2 to 134-8): archive cote 134-5, site page 77 (PDF page 78), typescript page 264.
% typescript. Draft compared: ps4.tex:3904-3965.
% AI-transcribed by Claude Opus 5.5 (claude-opus-5-5, adaptive thinking, effort high), Message Batches API, batch ps2 (msgbatch_01NVwr1Yuo8qJf5rykeGBXPK), 2026-11-12; not yet checked by eye.
\scanpage{134-5/77}
\mspage{264}

various homotopy notions in section 51\inserted{and 52}.  On the other hand,  it is clear
that the main ideas which are coming to fruition in part IV  all origina-%
ted \xout{in} during the awkward grinding process in part III  !

The  ``modelizing story''  so far has \xout{centered around}\inserted[typed]{turned out as} the interplay
of three main \xout{types}\inserted[typed]{sets} of notions.  One \xout{set} is made up with the \xout{``test-notions}
``\uline{test}-\uline{notions}'',  centering around the notion of a \uline{test}-\uline{category} ,  as
one giving rise to the most elementary type of ``\uline{modelizers}'' ,  namely the
so called ``\uline{elementary modelizers}'' $(A\hat{}\,,\underline{W}_A)$.  \xout{\unsure{This seemed}} This was develop-%
ped in part II  (while part I was concerned with the initial motivation
of the reflections,  namely stacks,  forgotten for the time being !).
The second set of notions concerns the so-called ``\uline{homotopy} \uline{notions}'' ,
developped at some length  in sections 51 to 55,  summarized in the dia-%
gram (2) above.  They constitute the main technical content of part III
of the notes,  with however one major shortcoming:  the relationship
between these notions,  and 0-connectedness notions  (20),  was only par-%
tially understood in part III, \inserted{namely} as a one-way relationship merely,  asso-%
ciating to \inserted{suitable} 0-connectedness notions in a category M,  a corresponding ho-%
motopy structure in M.  The third set of notions \xout{centers around the} may
be called \hand{``}\uline{asphericity} \uline{notions}'',  they center around the notions of
\uline{aspheric} \uline{objects} and \uline{aspheric} \uline{maps} (in a category endowed with a\xout{n} so-%
called \xout{a} \uline{asphericity} \uline{structure} ),  and more specifically around the
notion of \inserted{an} aspheric map in (Cat),  whose formal properties turn out to be
the \xout{main source} key for the development of a theory of asphericity struc-%
tures.  The first and the third set of notions  (namely test notions  and
asphericity notions)  depend on the choice of a ``\uline{basic localizer}'' \inserted{$\underline{W}$} in (Cat)\hand{\unsure{,}}
whereas the second set,  namely homotopy notions,  is ``absolute'' \deleted{namely}\hand{i.e.}
does not depend on any such choice,  nor on any knowledge of homotopy or
(co)homology invariants.

Whereas the test notions are essentially concerned with modelizers,
namely getting descriptions of the category of homotopy types (Hot)  in
terms of elementary modelizers $A\hat{}$ (as \inserted{being} $\underline{W}_A^{\ -1}A\hat{}$),  it appears that the
homotopy notions\hand{,} as well as the asphericity notions\hand{,} are independant of
any modelizing \inserted{notions and} assumptions.  In a deductive presentation of the theory,
the test notions would come last,  whereas they came first in these notes
-  as an illustration of the \inserted{general} fact that the deductive approach will present
things roughly in opposite order in which they have been discovered!  The
test notions,  as an outcome \inserted{of the attempt} to get a picture of modelizers,  have kept
acting as a constant \xout{source of} guideline in the whole reflection,  even
though technically speaking they are ``irrelevant'' for the development of
\margin{77}
